\documentclass[10pt,a4paper]{article}
\usepackage[utf8]{inputenc}
\usepackage[T1]{fontenc}
\renewcommand{\contentsname}{Contenido}
\renewcommand{\abstractname}{Resumen}
\emergencystretch=2em
\usepackage{geometry,booktabs,longtable,array,hyperref,xcolor,amsmath}
\geometry{margin=2cm}
\hypersetup{colorlinks=true,linkcolor=blue,urlcolor=blue}
\title{Trabajo de Título\\Etapa 2: análisis y diseño de Placement y Glance}
\author{Paolo Castañon Barrera \and Juan Claudio Cárdenas Uribe\\Profesor guía: Alejandro Mellado Gatica}
\date{2 de octubre de 2026}
\begin{document}
\maketitle
\begin{abstract}
Se presenta el análisis verificable de OpenStack Placement y Glance para una reimplementación acotada en Go. Una VM Ubuntu 24.04 dedicada quedó provisionada con Terraform y Ansible; DevStack, los smoke tests y el baseline original se ejecutaron con éxito. Los siete elementos de Sprint 2 suman 33 de 33 puntos aceptados con evidencias en el repositorio. La retroalimentación formal del profesor guía sigue pendiente de recibir.
\end{abstract}
\tableofcontents

\section{Introducción}
El proyecto compara compatibilidad funcional y, posteriormente, rendimiento y recursos de servicios originales Python frente a implementaciones Go. No se presupone una mejora de rendimiento. La planificación oficial es el libro \texttt{Planificacion\_Trabajo\_Titulo\_OpenStack\_Placement\_Glance.xlsx}, con fecha final de Sprint 2 el 2 de octubre de 2026.

\section{Estado respecto de Etapa 1}
El Excel contiene OE1--OE5, RF01--RF12, RNF01--RNF07 y backlog B01--B26. La Etapa 2 prioriza OE1 (análisis) y OE2 (entorno y baseline). En el repositorio inicial sólo existía un README; se incorporó una copia inalterada del libro y una matriz de requerimientos.

\section{Retroalimentación recibida}
No se encontró retroalimentación formal del profesor guía entre los insumos disponibles. \texttt{feedback-etapa1.md} deja el ítem pendiente para incorporación manual, sin inventar observaciones.

\section{Objetivos del Sprint 2}
OE1 exige arquitectura, APIs, dependencias y flujos críticos de Placement y Glance. OE2 exige OpenStack de un nodo reproducible, smoke tests y baseline original. RF03--RF10 describen implementación posterior; el diseño de esta etapa no los valida.

\section{Entorno experimental}
El host de trabajo CachyOS se mantuvo separado del laboratorio. Terraform creó la VM libvirt \texttt{tt-openstack-lab}, Ubuntu 24.04.5, 4 vCPU, 8 GiB asignados y disco de 55 GiB. Ansible verificó sistema, recursos, red y sudo, instaló MySQL y ejecutó DevStack como usuario \texttt{stack}. La auditoría de VM y los JSON de servicios están en \texttt{evidence/sprint2}.

\section{DevStack}
El instalador fija DevStack \texttt{stable/2026.2} al SHA \texttt{0fb9685c1bf22cd0b4d576d7a74f64195ad934fe} y aplica un parche versionado para inicializar Geneve antes de crear la red. Los servicios Keystone, Nova, Placement, Glance, Neutron, MySQL y RabbitMQ respondieron. \texttt{versions.lock} registra Placement \texttt{f4a89d3}, Glance \texttt{a3e5f49}, Nova \texttt{247ea5d}, Python 3.12.3, MySQL 8.0.46 y OpenStackClient 10.3.0. Ansible se repitió con éxito sin reinstalar. Secretos y logs crudos permanecen sólo en la VM.

\section{Arquitectura general}
\begin{center}
\fbox{\parbox{0.90\linewidth}{\centering Usuario / OpenStackClient $\rightarrow$ Keystone + Nova API\\
Nova API $\xrightarrow{RPC}$ Nova Scheduler $\xrightarrow{HTTP}$ Placement $\rightarrow$ DB\\
Nova Scheduler $\xrightarrow{RPC}$ Nova Compute $\xrightarrow{HTTP}$ Glance $\rightarrow$ image store\\
Nova / Compute $\xrightarrow{HTTP}$ Neutron}}
\end{center}
Las fuentes Mermaid completas están en \texttt{docs/sprint2/diagrams}. Las rutas HTTP utilizan tokens Keystone; las relaciones internas de Nova usan RPC.

\section{Diseño Placement}
El servicio desplegado Placement SHA \texttt{f4a89d3803cb23df22b8714569889ab39c918c34} expone handlers WSGI, microversiones, dominio y SQLAlchemy. ResourceProvider publica capacidad; Inventory cuantifica recursos; Trait y ResourceClass describen cualidades/tipos; Allocation y Consumer registran uso. \texttt{generation} controla concurrencia. El servidor anunció 1.0--1.39; el baseline probó providers, inventories, traits, classes, candidates, create/update/delete y 404. La traza de acceso mostró solicitudes de Nova y sus microversiones reales.

\section{Diseño Glance}
El servicio desplegado Glance SHA \texttt{a3e5f490b7a28b1161714c94acffe5308d24b567} separa API, dominio, DB de metadata y \texttt{glance\_store} para binarios. La configuración real usa \texttt{/opt/stack/data/glance/images/}. El baseline creó una imagen, cargó bytes, la consultó, modificó metadata, descargó bytes idénticos y la eliminó. Una petición sin token obtuvo HTTP 401; un boot desde CirrOS alcanzó ACTIVE.

\section{Contratos API}
\begin{center}\begin{tabular}{p{2.6cm}p{7.0cm}p{3.0cm}}\toprule
Servicio & Operaciones priorizadas & Estado de prueba\\\midrule
Placement & providers, inventories, traits, classes, allocations, candidates & Baseline ejecutado\\
Glance & POST/GET/PATCH/DELETE imágenes; PUT/GET file & Baseline ejecutado\\\bottomrule
\end{tabular}\end{center}
Placement anunció 1.0--1.39. La traza del boot registró GET candidatos y PUT allocations con 1.36, GET allocations con 1.28 y DELETE allocations con 1.0. Glance anunció Image API v2.18. Cada método/ruta, error y prioridad está en los inventarios API del repositorio.

\section{Modelos de datos}
\begin{center}\begin{tabular}{p{3.1cm}p{9.4cm}}\toprule
Dominio & Entidades verificadas en modelos SQLAlchemy\\\midrule
Placement & ResourceProvider, Inventory, ResourceClass, Trait, Allocation, Consumer; relaciones y generaciones\\
Glance & Image, ImageProperty, ImageLocation, ImageTag, ImageMember; metadata en DB y bytes en store\\\bottomrule
\end{tabular}\end{center}
Los diagramas ER indican claves y relaciones lógicas. En Placement no todas las relaciones ORM están declaradas como FK SQL.

\section{Arquitectura objetivo Go}
Se proponen procesos \texttt{cmd/placement} y \texttt{cmd/glance}; paquetes internos API, servicio, modelo y repositorio, más \texttt{glance/store}. Autenticación Keystone, configuración, middleware y logging son compartidos. Handlers negocian HTTP/microversión; servicios resuelven reglas; repositorios manejan transacciones; store transmite bytes.

\section{Estrategia de compatibilidad}
Cada petición se enviará al original y a Go con el mismo estado inicial. Se compararán status, JSON, tipos, campos, headers estables, errores y efectos persistidos. Nova y el cliente OpenStack serán verificadores adicionales. UUID, timestamps, orden no garantizado y request IDs se normalizarán sin ocultar diferencias semánticas. El baseline original y el smoke test se ejecutaron; los fixtures diferenciales contra Go corresponden a Sprint 3.

\section{Estrategia de pruebas}
Unitarios Go, contrato API, comparación diferencial, integración con Nova y end-to-end de VM forman la pirámide. En la fase posterior se medirán p50/p95/p99 de latencia, throughput, CPU, RAM y errores con carga controlada. No se presentan benchmarks ficticios.

\section{Decisiones técnicas}
Los ADR en \texttt{design-decisions.md} comparan schema original, propio e intermedio. Se propone schema propio Placement compatible por API y metadata Glance acotada con store filesystem único, ahora contrastado con el backend real. Auth Keystone se conservará. Las decisiones se revisarán durante las pruebas diferenciales.

\section{Estado del backlog}
\begin{center}\begin{tabular}{lll}\toprule
ID & Estado & Motivo principal\\\midrule
B01 & Hecho & VM dedicada auditada\\
B02 & Hecho & DevStack y endpoints verificados\\
B03 & Hecho & smoke y boot ACTIVE\\
B04 & Hecho & arquitectura y servicio contrastados\\
B05 & Hecho & alcance y microversiones medidas\\
B06 & Hecho & tráfico Nova a Placement capturado\\
B19 & Hecho & Image API, auth y store probados\\\bottomrule
\end{tabular}\end{center}
Los siete elementos S2 cumplen sus criterios con evidencia concreta. Avance formal: 7 de 7 elementos y 33 de 33 story points, es decir 100\% por la métrica de aceptación del Excel derivado.

\section{Riesgos}
Riesgos remanentes: dependencias secundarias no fijadas para reproducción bit a bit, cuerpos HTTP internos de Nova no capturados, cobertura de operaciones no críticas aún incompleta y retroalimentación formal del profesor no recibida. Los límites están en \texttt{BLOCKERS.md}.

\section{Conclusiones de la etapa}
OE1 y OE2 se contrastaron con DevStack real, APIs, base de datos, logs y un boot activo. Sprint 2 quedó aceptado con 33/33 puntos, sin anticipar implementación Go ni benchmarks. El Excel derivado, los scripts y las evidencias permiten iniciar Sprint 3.

\section{Próximos pasos Sprint 3}
Implementar B07 (providers/inventories), B08 (traits/classes) y B11 (persistencia), más B20/B25 de Glance según la planificación, usando los fixtures del baseline para pruebas diferenciales. Capturar cuerpos HTTP internos y ampliar errores/operaciones según lo requiera la compatibilidad. Incorporar retroalimentación formal del profesor cuando esté disponible.

\begin{thebibliography}{9}
\bibitem{devstack}\url{https://docs.openstack.org/devstack/latest/}
\bibitem{placement}\url{https://docs.openstack.org/api-ref/placement/}
\bibitem{glance}\url{https://docs.openstack.org/api-ref/image/v2/}
\bibitem{glancearch}\url{https://docs.openstack.org/glance/latest/contributor/architecture.html}
\bibitem{nova}\url{https://docs.openstack.org/nova/latest/admin/architecture}
\end{thebibliography}
\end{document}
